\documentclass[12pt,a4paper]{article}
\usepackage[utf8]{inputenc}
\usepackage[spanish,german]{babel}
\usepackage[T1]{fontenc}
\usepackage{geometry}
\geometry{a4paper, margin=2.5cm}
\usepackage{titlesec}
\usepackage{enumitem}
\usepackage{multicol}
\usepackage{csquotes}

\titleformat{\section}
{\Large\bfseries}
{}
{0em}
{}[\titlerule]

\titleformat{\subsection}
{\large\bfseries}
{}
{0em}
{}

\setlength{\parindent}{0pt}
\setlength{\parskip}{0.5em}

\begin{document}

\title{\textbf{DIALOG: Describir Objetos}}
\author{Spanisch A1/A2}
\date{}
\maketitle

\textbf{Tema:} Describir objetos \\
\textbf{Nivel:} A1 \\
\textbf{Complejidad:} Einfach \\
\textbf{Palabras:} 734

\vspace{1em}

\textit{Nota: Se recomienda revisar primero el vocabulario antes de leer el diálogo.}

\vspace{1em}

\section*{Diálogo}

\textbf{Camila:} Hola, Roberto. ¿Puedes ayudarme? Perdí mi celular y no lo encuentro.

\textbf{Roberto:} Claro, te ayudo. ¿Cómo es tu celular? ¿Puedes describirlo?

\textbf{Camila:} Sí, es un celular negro. Es pequeño y rectangular. Tiene una pantalla grande.

\textbf{Roberto:} ¿De qué marca es?

\textbf{Camila:} Es de marca Samsung. Es un modelo nuevo, del año pasado.

\textbf{Roberto:} ¿Tiene alguna funda o protector?

\textbf{Camila:} Sí, tiene una funda roja. La funda es de plástico y tiene un diseño de flores.

\textbf{Roberto:} Entiendo. ¿Y qué más tiene? ¿Tiene alguna marca o raya?

\textbf{Camila:} Sí, tiene una pequeña raya en la esquina superior derecha. Es muy pequeña, casi no se ve.

\textbf{Roberto:} Vale. ¿Y cómo es el tamaño? ¿Es grande o pequeño?

\textbf{Camila:} Es mediano. No es muy grande ni muy pequeño. Cabe perfectamente en mi bolso.

\textbf{Roberto:} ¿Y qué color tiene exactamente? ¿Es negro oscuro o negro claro?

\textbf{Camila:} Es negro oscuro, casi negro mate. No es brillante.

\textbf{Roberto:} Perfecto. Ahora dime, ¿dónde lo viste por última vez?

\textbf{Camila:} Lo vi por última vez en la mesa de la cocina. Estaba al lado de mi taza de café.

\textbf{Roberto:} ¿Y cómo era la taza? ¿Puedes describirla también?

\textbf{Camila:} La taza es blanca con dibujos azules. Es de cerámica y tiene un asa grande. Es redonda y no muy alta.

\textbf{Roberto:} ¿Y qué más había en la mesa?

\textbf{Camila:} Había un plato blanco, un vaso de agua y un libro. El libro es grande, tiene una portada azul y es bastante grueso.

\textbf{Roberto:} ¿Y el plato? ¿Cómo es?

\textbf{Camila:} El plato es blanco, redondo y plano. Es de porcelana. No tiene dibujos, es completamente blanco.

\textbf{Roberto:} Entiendo. Oye, ¿recuerdas si había algo más en la mesa? ¿Algún objeto pequeño?

\textbf{Camila:} Sí, había unas llaves. Las llaves son plateadas y tienen un llavero rojo. El llavero tiene forma de corazón.

\textbf{Roberto:} Perfecto. Ahora voy a buscar en la cocina. ¿Puedes describir también la mesa?

\textbf{Camila:} Claro. La mesa es de madera, de color marrón claro. Es rectangular y tiene cuatro patas. Es bastante grande.

\textbf{Roberto:} ¿Y qué hay encima de la mesa normalmente?

\textbf{Camila:} Normalmente hay un florero con flores. El florero es de vidrio, transparente y alto. Las flores son de diferentes colores.

\textbf{Roberto:} Muy bien. Déjame buscar. Oye, ¿has revisado tu bolso?

\textbf{Camila:} Sí, lo revisé. Mi bolso es negro, de cuero. Es mediano, tiene dos asas y un cierre de metal. Pero no está allí.

\textbf{Roberto:} ¿Y cómo es el cierre? ¿Es dorado o plateado?

\textbf{Camila:} Es plateado, de metal brillante. Tiene un diseño simple.

\textbf{Roberto:} Entiendo. Déjame buscar en la cocina. ¿Puedes describir también la silla que está cerca de la mesa?

\textbf{Camila:} La silla es de madera, del mismo color que la mesa. Tiene un respaldo alto y es bastante cómoda. Tiene cuatro patas también.

\textbf{Roberto:} Perfecto. Voy a buscar ahora. Espera un momento.

\textbf{Camila:} Gracias, Roberto. Te agradezco mucho tu ayuda.

\textbf{Roberto:} De nada. Es un placer ayudarte. Oye, ¿encontraste algo en tu bolso que pueda ayudarnos?

\textbf{Camila:} Déjame ver... Tengo mi billetera, que es negra y pequeña. También tengo unas gafas de sol, que son negras con montura dorada. Pero no está mi celular.

\textbf{Roberto:} No te preocupes. Vamos a encontrarlo. ¿Quieres que busquemos juntos?

\textbf{Camila:} Sí, por favor. Eso sería muy útil.

\textbf{Roberto:} Perfecto. Empecemos por la cocina.

\newpage

\section*{Vocabulario}

\begin{multicols}{2}
\begin{itemize}[leftmargin=*]
    \item \textbf{objeto} - Gegenstand, Objekt
    \item \textbf{describir} - beschreiben
    \item \textbf{celular} - Handy, Mobiltelefon
    \item \textbf{rectangular} - rechteckig
    \item \textbf{pantalla} - Bildschirm, Display
    \item \textbf{marca} - Marke
    \item \textbf{modelo} - Modell
    \item \textbf{funda} - Hülle, Schutzhülle
    \item \textbf{protector} - Schutz, Schutzhülle
    \item \textbf{plástico} - Plastik
    \item \textbf{diseño} - Design, Muster
    \item \textbf{raya} - Kratzer, Ritz
    \item \textbf{esquina} - Ecke
    \item \textbf{superior} - oberer/obere/oberes
    \item \textbf{mediano} - mittelgroß
    \item \textbf{bolso} - Handtasche, Tasche
    \item \textbf{mate} - matt (nicht glänzend)
    \item \textbf{brillante} - glänzend, glitzernd
    \item \textbf{taza} - Tasse
    \item \textbf{cerámica} - Keramik
    \item \textbf{asa} - Henkel, Griff
    \item \textbf{redondo} - rund
    \item \textbf{plato} - Teller
    \item \textbf{porcelana} - Porzellan
    \item \textbf{llaves} - Schlüssel
    \item \textbf{plateado} - silberfarben
    \item \textbf{llavero} - Schlüsselanhänger
    \item \textbf{corazón} - Herz
    \item \textbf{madera} - Holz
    \item \textbf{marrón} - braun
    \item \textbf{patas} - Beine (von Möbeln)
    \item \textbf{florero} - Vase
    \item \textbf{vidrio} - Glas
    \item \textbf{transparente} - transparent, durchsichtig
    \item \textbf{cuero} - Leder
    \item \textbf{cierre} - Verschluss
    \item \textbf{metal} - Metall
    \item \textbf{dorado} - goldfarben
    \item \textbf{respaldo} - Rückenlehne
    \item \textbf{billetera} - Geldbörse, Portemonnaie
    \item \textbf{gafas de sol} - Sonnenbrille
    \item \textbf{montura} - Brillengestell
\end{itemize}
\end{multicols}

\newpage

\section*{Preguntas de Comprensión}

\begin{enumerate}[leftmargin=*]
    \item ¿Qué perdió Camila?
    \item ¿De qué color es el celular de Camila?
    \item ¿De qué marca es el celular?
    \item ¿Qué tiene el celular en la esquina superior derecha?
    \item ¿Cómo es el tamaño del celular?
    \item ¿Dónde vio Camila el celular por última vez?
    \item ¿Cómo es la taza de café que menciona Camila?
    \item ¿Qué más había en la mesa además del celular y la taza?
    \item ¿Cómo es el libro que estaba en la mesa?
    \item ¿Cómo es el plato?
    \item ¿Qué objeto pequeño había en la mesa?
    \item ¿Cómo es el llavero de las llaves?
    \item ¿De qué material es la mesa?
    \item ¿Qué hay normalmente encima de la mesa?
    \item ¿Cómo es el florero?
    \item ¿Cómo es el bolso de Camila?
    \item ¿De qué color es el cierre del bolso?
    \item ¿Cómo es la silla que está cerca de la mesa?
    \item ¿Qué objetos tiene Camila en su bolso?
    \item ¿Cómo son las gafas de sol de Camila?
\end{enumerate}

\newpage

\section*{Verbo Irregular: VER (sehen)}

\textit{Nota: \enquote{vosotros/as} se usa solo en España, no en español sudamericano. Se incluye aquí para completitud.}

\begin{multicols}{2}

\subsection*{Presente}
\begin{itemize}[leftmargin=*]
    \item yo - veo
    \item tú - ves
    \item él/ella/usted - ve
    \item nosotros/as - vemos
    \item vosotros/as - veis
    \item ustedes - ven
    \item ellos/ellas - ven
\end{itemize}

\subsection*{Pretérito Indefinido}
\begin{itemize}[leftmargin=*]
    \item yo - vi
    \item tú - viste
    \item él/ella/usted - vio
    \item nosotros/as - vimos
    \item vosotros/as - visteis
    \item ustedes - vieron
    \item ellos/ellas - vieron
\end{itemize}

\subsection*{Pretérito Imperfecto}
\begin{itemize}[leftmargin=*]
    \item yo - veía
    \item tú - veías
    \item él/ella/usted - veía
    \item nosotros/as - veíamos
    \item vosotros/as - veíais
    \item ustedes - veían
    \item ellos/ellas - veían
\end{itemize}

\columnbreak

\subsection*{Futuro Simple}
\begin{itemize}[leftmargin=*]
    \item yo - veré
    \item tú - verás
    \item él/ella/usted - verá
    \item nosotros/as - veremos
    \item vosotros/as - veréis
    \item ustedes - verán
    \item ellos/ellas - verán
\end{itemize}

\subsection*{Pretérito Perfecto}
\begin{itemize}[leftmargin=*]
    \item yo - he visto
    \item tú - has visto
    \item él/ella/usted - ha visto
    \item nosotros/as - hemos visto
    \item vosotros/as - habéis visto
    \item ustedes - han visto
    \item ellos/ellas - han visto
\end{itemize}

\end{multicols}

\newpage

\section*{Explicación Gramatical}

\textbf{Adjektive zur Beschreibung von Objekten}

Im Spanischen werden Adjektive verwendet, um Objekte zu beschreiben. Wichtig ist die Übereinstimmung (Kongruenz) zwischen Substantiv und Adjektiv in Geschlecht (maskulin/feminin) und Zahl (Singular/Plural).

\textbf{Adjektivstellung:}
\begin{itemize}[leftmargin=*]
    \item \textbf{Nach dem Substantiv} (meistens): \enquote{un celular negro} (ein schwarzes Handy)
    \item \textbf{Vor dem Substantiv} (bei subjektiven/emotionalen Adjektiven): \enquote{un buen libro} (ein gutes Buch)
    \item \textbf{Beide Positionen möglich} (Bedeutung kann sich ändern): \enquote{un hombre grande} (ein großer Mann) vs. \enquote{un gran hombre} (ein großer/berühmter Mann)
\end{itemize}

\textbf{Kongruenz:}
\begin{itemize}[leftmargin=*]
    \item Maskulin Singular: \enquote{el celular negro} (das schwarze Handy)
    \item Feminin Singular: \enquote{la taza blanca} (die weiße Tasse)
    \item Maskulin Plural: \enquote{los celulares negros} (die schwarzen Handys)
    \item Feminin Plural: \enquote{las tazas blancas} (die weißen Tassen)
\end{itemize}

\textbf{Beispiele aus dem Text:}
\begin{itemize}[leftmargin=*]
    \item \enquote{es un celular negro} -- es ist ein schwarzes Handy (maskulin, singular)
    \item \enquote{es pequeño y rectangular} -- es ist klein und rechteckig (maskulin, singular)
    \item \enquote{tiene una pantalla grande} -- es hat einen großen Bildschirm (feminin, singular)
    \item \enquote{una funda roja} -- eine rote Hülle (feminin, singular)
    \item \enquote{es de plástico} -- es ist aus Plastik (Material)
    \item \enquote{la taza es blanca} -- die Tasse ist weiß (feminin, singular)
    \item \enquote{con dibujos azules} -- mit blauen Zeichnungen (maskulin, plural)
    \item \enquote{es de cerámica} -- es ist aus Keramik (Material)
    \item \enquote{el plato es blanco} -- der Teller ist weiß (maskulin, singular)
    \item \enquote{las llaves son plateadas} -- die Schlüssel sind silberfarben (feminin, plural)
    \item \enquote{la mesa es de madera} -- der Tisch ist aus Holz (Material)
    \item \enquote{de color marrón claro} -- von hellbrauner Farbe
    \item \enquote{el florero es de vidrio} -- die Vase ist aus Glas (Material)
    \item \enquote{es transparente} -- es ist transparent (feminin, singular)
    \item \enquote{mi bolso es negro} -- meine Handtasche ist schwarz (maskulin, singular)
    \item \enquote{de cuero} -- aus Leder (Material)
\end{itemize}

\textbf{Wichtige Adjektive für Objektbeschreibungen:}
\begin{itemize}[leftmargin=*]
    \item \textbf{Größe}: grande (groß), pequeño (klein), mediano (mittelgroß), alto (hoch), bajo (niedrig)
    \item \textbf{Form}: redondo (rund), rectangular (rechteckig), cuadrado (quadratisch), largo (lang), corto (kurz)
    \item \textbf{Farbe}: negro (schwarz), blanco (weiß), rojo (rot), azul (blau), verde (grün), etc.
    \item \textbf{Material}: de plástico (aus Plastik), de madera (aus Holz), de metal (aus Metall), de vidrio (aus Glas), de cuero (aus Leder)
    \item \textbf{Zustand}: nuevo (neu), viejo (alt), limpio (sauber), sucio (schmutzig), roto (kaputt)
\end{itemize}

\textbf{Hinweis}: Bei Materialangaben verwendet man \enquote{de + Material} (z.B. \enquote{de madera}, \enquote{de plástico}). Das Adjektiv steht meistens nach dem Substantiv.

\end{document}
